\chapter{Chapitre de test du préambule}
\label{chap:smoke}

\begin{objectifs}
  \item Vérifier que tous les environnements compilent.
  \item Vérifier les caractères Unicode dans les extraits : ⊤, ⊥, →, ─, é.
\end{objectifs}

\prerequis{aucun.}

\section{Titre avec \code{snake_case} et \code{s0=⊤}}
Cas : \code{a && b}, \code{x \% 2}, \code{\#[test]}, \code{\{ ...q \}}, \code{$x^2$}, $n \in \code{s_0}$,
\code{a~b}, \code{✓ 1 file(s) no issues found.}, \code{règle → ⊥}, \code{\ldots},
\file{src/ir/expr.rs}, \code{new_state.leq(&state)}, fin.

Texte avec \code{snake_case_ident} et \file{src/ir/expr.rs}, une règle \regle{infinite-loop},
une relation \relation{writers}, l'\adr{27}, l'issue \issue{63}, et un \terme{treillis}.
Prose Unicode : ⊤ ⊥ ⊓ ⊔ → ⇒ ≤ ∀ ∞. Math : $x \lleq y \join z$, $\gam(\Top) = \Val$, $f^{\widen}$.

\begin{definition}{Treillis}{treillis}
  Un ensemble partiellement ordonné $(L, \lleq)$ dont toute paire admet une borne supérieure $\join$ et une borne inférieure $\meet$.
\end{definition}

\begin{react}[batching]
  Plusieurs appels à \code{setState} dans le même tick sont fusionnés.
\end{react}

\begin{piege}
  Un faux négatif est interdit.
\end{piege}

\begin{exemple}{Compteur}{compteur}
Renvoi vers la \cref{def:treillis} et l'\cref{ex:compteur} ; deux points : ici.
\begin{lstlisting}[language=TypeScript,caption={Un compteur — é ⊤ → …}]
// … points de suspension …
function Counter() {
  const [n, setN] = useState(0);          // état local
  useEffect(() => { setN(n + 1); }, [n]); // ⊤ → boucle
  return <button onClick={() => setN(n + 1)}>{n}</button>;
}
\end{lstlisting}
\end{exemple}

\rustsnippet{src/ir/expr.rs}{1}{25}{en-tête du module des expressions}

\begin{sortie}
error[infinite-loop]: useEffect writes `n` and depends on it — é ⊤ → ─────
\end{sortie}

\ifhastikz
\begin{figure}[H]\centering
\begin{tikzpicture}[node distance=8mm]
  \node[bloc] (a) {parse (oxc)};
  \node[bloc,right=of a] (b) {lowering};
  \node[bloc,right=of b] (c) {IR};
  \node[bloc,right=of c] (d) {engine};
  \node[bloc,right=of d] (e) {rules};
  \draw[fleche] (a) -- (b); \draw[fleche] (b) -- (c); \draw[flechemust] (c) -- (d); \draw[flechemay] (d) -- (e);
\end{tikzpicture}
\caption{Pipeline (TikZ disponible).}
\end{figure}
\else
\emph{TikZ indisponible : figure omise.}
\fi

\begin{exemple}{}{}
  Exemple sans titre ni label.
\end{exemple}

\begin{algorithm}[H]
\KwIn{un CFG $G$, un état initial $s_0$}
\KwOut{une table $S$ d'états abstraits}
$S[\mathit{entry}] \gets s_0$\;
\While{la worklist n'est pas vide}{
  $b \gets$ pop\;
  $s \gets \bigsqcup_{p \in \pred(b)} S[p]$\;
  \If{$b$ est une tête de boucle}{$s \gets S[b] \widen s$\;}
}
\caption{Itération de point fixe (test).}
\end{algorithm}

\begin{exercice}{Un exercice}{un}
  Question ?
\end{exercice}
\begin{solution}
  Réponse~\cite{CousotCousot77}.
\end{solution}

\begin{resume}
Tout compile.
\end{resume}
